\begin{exsol}{7}{Varian Ch.~35, Review Question~1}{687}{Foundations}{B1}
\asked Judge the claim that clearly defining property rights usually removes the problem of
externalities, saying what such a delineation achieves and what it leaves open.

\method Externalities as missing or ill-defined property rights (Varian \S35.1, pp.~667--670);
quasilinear preferences and the Coase theorem (Varian \S35.2, pp.~670--672); \MR's value
maximization principle and Coase theorem with their hypotheses (pp.~35--39); bargaining costs,
private information and many parties (\MR{} Ch.~5, pp.~140--149).

\solution
\textbf{Verdict: true as a statement about efficiency in small-numbers, low-cost bargaining; false
as a general policy claim.} Varian's own answer (p.~A29) is ``True'' with a distributional
caveat. We agree with it and make the qualifications explicit.

\emph{(a) What delineation fixes.} An externality is an effect of one agent's action on another's
utility or production that is not priced. Varian's diagnosis is that ``the practical problems with
externalities generally arise because of poorly defined property rights'' (p.~670): if the smoker
believes he may smoke and the roommate believes she is entitled to clean air, nobody can buy or
sell the right, and the level of smoke is set by whoever acts. Once the right is assigned and
enforceable, it becomes an endowment that can be traded like money or bread. Starting from that
endowment, voluntary trade moves the parties to the contract curve, where
\[
  MRS^A_{\text{smoke},\,\text{money}} \;=\; MRS^B_{\text{smoke},\,\text{money}},
\]
so the level of the externality is Pareto efficient \emph{whichever} party holds the right
(Varian p.~669, points $X$ and $X'$ of Figure~35.1).

\emph{(b) What delineation does not fix.}
\begin{enumerate}[label=(\roman*)]
  \item \emph{Distribution.} Assigning the right is assigning an endowment. The holder is paid
  to give it up rather than paying to obtain it, so the final money holdings differ across
  assignments (Varian p.~A29; item~8 computes the difference).

  \item \emph{The efficient level itself, when there are wealth effects.} Without quasilinearity
  the efficient amount of the externality depends on who holds the right (Varian p.~670).
  Example: a plant gains $G = 40$ from dumping into a river; a resident with wealth $w = 100$ has
  $u = \ln x + a\,y$, $y = 1$ if the river is clean, with $a = \ln 1.5$. Then
  \[
    \text{WTP} = w(1 - e^{-a}) = 100\big(1 - \tfrac{2}{3}\big) = 33.33, \qquad
    \text{WTA} = w(e^{a} - 1) = 100(1.5 - 1) = 50 .
  \]
  If the plant holds the right, the resident can offer at most $33.33 < 40$: the river is
  polluted. If the resident holds it, the plant would have to pay at least $50 > 40$: the river
  stays clean. Both outcomes are efficient given their endowment, yet the \emph{real} outcome
  differs. Delineation still delivers efficiency, but ``the'' efficient level no longer exists
  independently of the assignment (this is the gap between value maximization and efficiency
  under wealth effects, \MR{} n.~10, p.~37).

  \item \emph{Transaction costs.} The Coase theorem assumes that the parties ``bargain to an
  efficient agreement'' (\MR{} p.~38). That fails when
  \begin{itemize}
    \item there are many parties (air pollution over a city): each victim free rides on the
    others' payments, and the chance that anyone volunteers stays bounded away from one even when
    the project is surely value maximizing (\MR{} pp.~145--147);
    \item valuations are private: no bargaining procedure may reach the efficient trade
    (\MR{} pp.~141--143);
    \item measuring and enforcing the right is costly (diffuse emissions, noise at 3~a.m.), or
    bargaining itself absorbs resources (\MR{} pp.~147--149);
    \item parties are wealth constrained and cannot pay the compensation that efficiency
    requires (\MR{} p.~39).
  \end{itemize}
  In these cases, rights alone leave an inefficiency and other tools are needed: Pigouvian taxes,
  quotas, tradable permits, or mergers that internalize the effect (Varian \S35.3 and \S35.6).
\end{enumerate}

\interp The question is a compressed statement of the Coase theorem and its limits. Under its
hypotheses, assigning rights decides only the transfers; that is why \MR{} turn the theorem into
a research programme: look for the transaction costs that stop the parties from bargaining to the
efficient $y$, and explain organizations as ways to reduce them (p.~38). ``Property rights solve
externalities'' is the benchmark, not the conclusion.

\verify With $a = \ln 1.5$: $e^{-a} = 2/3$, so WTP $=100/3 \approx 33.33$; $e^{a} - 1 = 0.5$, so WTA
$= 50$; and $33.33 < 40 < 50$, so the efficient outcome flips with the assignment. Utility checks:
$\ln(100 - 100/3) + \ln 1.5 = \ln 100$ and $\ln(150) = \ln 100 + \ln 1.5$ (\texttt{check\_C.py}).

\answer{True only in the Coasean benchmark (well-defined, enforceable rights; few parties;
costless bargaining with full information): rights turn the externality into a tradable good and
trade reaches an efficient level. Even then the assignment changes the distribution, and with
wealth effects it changes the efficient level itself (WTP $33.33$ vs.\ WTA $50$ around a gain of
$40$). With many parties, private information or costly enforcement, rights alone do not remove
the inefficiency.}
\end{exsol}
